\documentclass[11pt]{article}
\usepackage[margin=1.1in]{geometry}
\usepackage{amsmath,amssymb,amsthm}
\usepackage{booktabs}
\usepackage{hyperref}

\newtheorem{lemma}{Lemma}
\newtheorem{remark}{Remark}

\newcommand{\bA}{A}
\newcommand{\bB}{B}
\newcommand{\bC}{C}
\newcommand{\Chat}{\hat C}
\newcommand{\khat}{\hat k}

\title{Finiteness of the longitudinal integral and subtraction of the
$y'=y$ point in the two--Wilson-line term}
\author{}
\date{\today}

\begin{document}
\maketitle

\section{The expression and the apparent divergence}

We study the term
\begin{align}
\sigma(k) \;=\; P\int_{k^{+}/V}^{1}d\xi\,\frac{1-\xi}{\xi}
\int_{x,y,x',y'} e^{-ik\cdot(y'-y)/\xi}\;
\mathcal K(x,y;x',y')\;
\big[U(y')-U(x')\big]^{ce'}\big[U(y)-U(x)\big]^{ce}\,
\rho^{e'}(x')\,\rho^{e}(x),
\label{eq:sigma}
\end{align}
with the prefactor $P=\frac{1}{(2\pi)^{3}}\frac{g^{4}}{8\pi^{4}}\frac{N_{c}}{k^{+}}$,
the longitudinal cutoff $V$, and the kernel
\begin{align}
\mathcal K
=\frac{(x-y)\cdot(x'-y')}{(x-y)^{2}(x'-y')^{2}}
-2\,\frac{(x-y)\cdot(y'-y)\,(x'-y')\cdot(y'-y)}
       {(x-y)^{2}(x'-y')^{2}(y'-y)^{2}}.
\label{eq:kernel}
\end{align}
Setting the phase to unity, the $\xi$ integral gives
$\int_{k^+/V}d\xi\,(1-\xi)/\xi=\ln(V/k^+)+\mathcal O(1)$, so
\eqref{eq:sigma} appears to be logarithmically divergent as $V\to\infty$,
and on the slice $y'=y$ (where the phase \emph{is} unity for all $\xi$)
the integrand genuinely is.  The purpose of this note is threefold:
(i) to show that for $k\neq0$ the integral \eqref{eq:sigma} nevertheless
has a finite $V\to\infty$ limit; (ii) to clarify in what sense the point
$y'=y$ is and is not divergent; and (iii) to construct an explicit
subtraction of the $y'=y$ point, with the counterterm in closed form,
so that the subtracted integrand is regular at $y'=y$ even at
$V=\infty$ --- the form suited for numerical evaluation.  All statements
are verified numerically in Sec.~\ref{sec:numerics}.

\section{Structure of the kernel}
\label{sec:kernel}

Write $\bA=x-y$, $\bB=x'-y'$, $\bC=y'-y$, $\Chat=\bC/|\bC|$.  The kernel
factorizes through the reflection tensor
\begin{align}
\mathcal K=\frac{\bA_i}{\bA^2}\,S_{ij}(\Chat)\,\frac{\bB_j}{\bB^2},
\qquad
S_{ij}(\hat n)=\delta_{ij}-2\hat n_i\hat n_j .
\label{eq:factor}
\end{align}
Three properties of this structure control everything below.

\paragraph{(i) Boundedness at $\bC\to0$.}  The $(y'-y)^2$ in the
denominator of the second term of \eqref{eq:kernel} is cancelled by the
two explicit factors of $(y'-y)$ in its numerator: $\mathcal K$ is a
$0/0$ of homogeneity degree zero in $\bC$, direction-dependent but
bounded.  Since $S$ is an orthogonal (reflection) matrix,
\begin{align}
|\mathcal K|\le \frac{1}{|\bA|\,|\bB|}\qquad\text{for all }y'.
\label{eq:bound}
\end{align}
There is no power singularity at $y'=y$; the only divergence anywhere
in \eqref{eq:sigma} is the longitudinal one.

\paragraph{(ii) Pure second harmonic.}  In the angle $\theta$ of
$\hat n$, $S_{11}=-\cos2\theta$, $S_{12}=S_{21}=-\sin2\theta$,
$S_{22}=\cos2\theta$: $S$ is a pure $m=\pm2$ harmonic
(equivalently $S_{ij}(\Chat)/\bC^2=\partial_i\partial_j\ln|\bC|$ away
from $\bC=0$), so $\oint d\theta\,S_{ij}(\hat n)=0$.

\paragraph{(iii) Exact angular cancellation.}  The angular average of
the \emph{full} kernel over the direction of $\bC$ vanishes identically,
not only as $|\bC|\to0$:

\begin{lemma}\label{lem:loop}
Let $\varphi$ be harmonic on the disk $|y'-y|<r_{\max}$.  Then for every
$r<r_{\max}$ and every fixed vector $\bA$,
\[
\oint_0^{2\pi} d\theta\; \bA_i\,S_{ij}(\Chat)\,
\partial_j\varphi\big(y+r\Chat\big)=0 ,
\qquad \Chat=(\cos\theta,\sin\theta).
\]
\end{lemma}

\begin{proof}
Use complex notation $z=\bC_1+i\bC_2=re^{i\theta}$ and represent a
vector $v$ by $v_1+iv_2$.  A reflection acts as
$S(\hat n)v=-e^{2i\theta}\bar v$.  Writing
$\varphi=\mathrm{Re}\sum_{m\ge0}a_m z^m$ (regular inside the disk), the
gradient is $h_1+ih_2=2\partial_{\bar z}\varphi=\sum_m m\bar a_m\bar
z^{\,m-1}$, hence
$S h=-e^{2i\theta}\overline{(h_1+ih_2)}
=-\sum_m m a_m r^{m-1}e^{i(m+1)\theta}$.
Contracting with $\bA$ gives a sum of terms $\propto e^{i(m+1)\theta}$
with $m\ge0$, each of which integrates to zero over $\theta$.
\end{proof}

Since $\bB_j/\bB^2=\partial_{y'_j}\big(-\ln|x'-y'|\big)$ and
$-\ln|x'-y'|$ is harmonic in $y'$ for $y'\neq x'$, Lemma~\ref{lem:loop}
applied to \eqref{eq:factor} gives
\begin{align}
\oint d\theta_{\Chat}\,\mathcal K\big(y'=y+r\Chat\big)=0
\qquad\text{for all } r<|x'-y| .
\label{eq:loopzero}
\end{align}
(The numerics of Sec.~\ref{sec:numerics} confirm this at machine
precision at all radii tested.)  In particular, the angular average of
the second term in \eqref{eq:kernel} equals the first exactly: the
second term is precisely the structure needed for the soft region to
decouple.

\section{Absence of $\ln V$ at fixed $k\neq0$}
\label{sec:finite}

Isolate the $\xi$ integral at fixed $y'$, with $a=k\cdot(y'-y)$ and
$\varepsilon=k^+/V$:
\begin{align}
F_\varepsilon(a)=\int_{\varepsilon}^{1}\frac{d\xi}{\xi}\,(1-\xi)\,
e^{-ia/\xi}
=\int_{1}^{1/\varepsilon}dt\left(\frac1t-\frac1{t^{2}}\right)e^{-iat},
\label{eq:F}
\end{align}
given in closed form in Appendix~\ref{app:F}.  Its two essential
properties are
\begin{align}
F_\varepsilon(0)=\ln\frac1\varepsilon-1+\varepsilon ,
\qquad
\big|F_\varepsilon(a)\big|\;\le\;
\ln\frac{1}{\max(|a|,\varepsilon)}+c_0\quad(|a|\le1),
\qquad
\big|F_\varepsilon(a)\big|\le\frac{c_1}{|a|}\quad(|a|\ge1),
\label{eq:Fbounds}
\end{align}
with $\varepsilon$-independent constants $c_{0,1}$: the oscillation of
the phase sets in at $\xi\sim|a|$ and decoheres the region $\xi<|a|$, so
the effective lower limit of the $\xi$ integral is
$\max(\varepsilon,|a|)$.  The cutoff $V$ therefore survives only where
$|k\cdot(y'-y)|\le k^+/V$.

Now estimate the $y'$ integral near the dangerous set.  Note that the
phase is frozen not only at the point $y'=y$ but on the whole line
$k\cdot(y'-y)=0$; the frozen region $|k\cdot\bC|\le\varepsilon$ is a
\emph{strip} of width $2\varepsilon/|k|$ around that line.  Within a
neighborhood $|\bC|<\delta$ of $y'=y$, using the uniform bound
\eqref{eq:bound} on the kernel and boundedness of the Wilson-line
factors:
\begin{itemize}
\item[\emph{(a)}] \emph{Frozen strip}, $|k\cdot\bC|\le\varepsilon$:
contribution bounded by
$\frac{2\varepsilon}{|k|}\cdot2\delta\cdot
M\,\ln\frac1\varepsilon\;\sim\;\frac{\ln V}{V}\;\xrightarrow{V\to\infty}\;0$.
The $\ln V$ in the integrand is real, but it sits on a region whose
measure shrinks like $1/V$.
\item[\emph{(b)}] \emph{Active region}, $|k\cdot\bC|>\varepsilon$: here
$|F_\varepsilon|\le\ln\frac{1}{|k\cdot\bC|}+c_0$ \emph{independently of}
$V$, and $\int_{|\bC|<\delta}d^2\bC\,\ln\frac{1}{|k\cdot\bC|}<\infty$
because a logarithm is integrable in one dimension.
\end{itemize}
Since $\ln(1/\varepsilon)\le\ln(1/|a|)$ exactly where the cutoff bites,
$\ln(1/|a|)+c_0$ is a $V$-independent integrable majorant and dominated
convergence applies:
\begin{align}
\boxed{\;\lim_{V\to\infty}\sigma(k)\ \text{exists and is finite for
every fixed } k\neq0 .\;}
\end{align}
No $\ln V$ remains.  The would-be rapidity logarithm has been traded for
the \emph{transverse} logarithm $\ln\frac{1}{|k||y'-y|}$, integrable
over $y'$.  By property (iii) of Sec.~\ref{sec:kernel} even this
logarithm is softer than it looks: its coefficient at $y'\to y$ is a
pure second harmonic in $\Chat$ and drops after the angular part of the
$y'$ integration, so the first surviving contribution from $y'\approx y$
comes from the $\mathcal O(|\bC|)$ variation of $U(y')$ and of
$\bB_j/\bB^2$.

\begin{remark}
The limit is not uniform in $k$.  At $k=0$ the phase is frozen on the
entire plane, the strip argument fails, and the $\ln V$ divergence is
genuine, with coefficient
$\propto\int_{x,y,x',y'}\mathcal K\,[U\cdots][U\cdots]\rho\rho$
(nonzero in general: Eq.~\eqref{eq:loopzero} kills only the
$y'\approx y$ region, not the whole integral, once the $y'$ dependence
of the Wilson-line factors is kept).  Correspondingly the finite answer
at $k\neq0$ contains $\ln(1/|k|\cdot\text{distances})$: the divergence
has been pushed to the corner $k\to0$, not removed from the theory.
\end{remark}

\begin{remark}
The point $y'=x'$, where the first term of \eqref{eq:kernel} behaves as
$1/|x'-y'|$, is protected separately: it is multiplied by
$[U(y')-U(x')]$, which vanishes linearly there.  Likewise $x=y$ is
protected by $[U(y)-U(x)]$.
\end{remark}

\section{The slice $y'=y$: divergent integrand, convergent integral}
\label{sec:slice}

On the slice $y'=y$ the phase equals $1$ for all $\xi$ and the integrand
carries the factor $F_\varepsilon(0)=\ln(V/k^+)-1+k^+/V\to\infty$: the
\emph{integrand} has no finite $V\to\infty$ limit at that point.  This
does not make the \emph{integral} divergent: $\{y'=y\}$ (and the line
$k\cdot\bC=0$) has zero two-dimensional measure, and by the estimates of
Sec.~\ref{sec:finite} the contribution of a neighborhood $|\bC|<\epsilon$
is bounded by $\epsilon^{2}\ln(1/\epsilon)+\mathrm{const}\cdot\epsilon$,
\emph{uniformly in $V$}.  As a function of $y'$ the $\xi$-integrated
integrand is a $\ln(V/k^+)$ plateau of width $\sim k^+/(|k|V)$ sitting
on the $V$-independent profile $\ln\frac{1}{|k||y'-y|}$; the plateau
contributes $\sim V^{-1}\ln V\to0$.  (Compare
$\int_0^1dy\,\ln\frac{1}{y+\epsilon}$: the integrand diverges at $y=0$
as $\epsilon\to0$, the integral does not.)  The pointwise limit
$V\to\infty$ and the $y'$ integration do not commute at $y'=y$;
exchanging them for the \emph{value} of the integral is justified by
dominated convergence.

\section{Subtraction of the $y'=y$ point}
\label{sec:subtraction}

For numerics one wants the integrand itself, not only the integral, to
be finite at $y'=y$ at $V=\infty$.  Expand the slowly varying factors
about $y'=y$, keeping the phase and the angular tensor exact:
\begin{align}
\mathcal K\,\big[U(y')-U(x')\big]
=\frac{\bA_i}{\bA^2}\,S_{ij}(\Chat)\,\frac{(x'-y)_j}{(x'-y)^{2}}
\,\big[U(y)-U(x')\big]\;+\;\mathcal O(|\bC|) .
\label{eq:expand}
\end{align}
Define the counterterm from the leading term, damped at large $|\bC|$ by
$\chi(|\bC|)$ (below, $\chi=e^{-\bC^2/r_0^2}$):
\begin{align}
\sigma_{\rm CT}=P\int_{k^+/V}^{1}d\xi\,\frac{1-\xi}{\xi}
\int_{x,y,x',y'} e^{-ik\cdot\bC/\xi}\,\chi(|\bC|)\,
\frac{(x-y)_i}{(x-y)^2}\,S_{ij}(\Chat)\,\frac{(x'-y)_j}{(x'-y)^2}\,
\big[U(y)-U(x')\big]^{ce'}\big[U(y)-U(x)\big]^{ce}\rho^{e'}\rho^{e},
\label{eq:CT}
\end{align}
and $\sigma_{\rm reg}=\sigma-\sigma_{\rm CT}$ \emph{at integrand level}.
By \eqref{eq:expand} the subtracted integrand differs from zero at
$y'\to y$ by $\mathcal O(|\bC|)$, so after the $\xi$ integration the
subtracted $y'$-density vanishes like $|\bC|\ln(1/|\bC|)$ at $y'\to y$:
the pole is removed pointwise and in $\sigma_{\rm reg}$ the cutoff may
be dropped under the integral.

\paragraph{The counterterm in closed form.}  The $y'$ integral in
\eqref{eq:CT} is a pure $J_2$ Hankel transform.  Using
$\int_0^{2\pi}d\theta\,e^{i(n\theta-x\cos\theta)}=2\pi(-i)^nJ_n(x)$ with
$n=2$,
\begin{align}
\int d^2\bC\;e^{-iq\cdot\bC}\,S_{ij}(\Chat)\,\chi(|\bC|)
=-2\pi\,S_{ij}(\hat q)\,T(|q|),
\qquad
T(q)=\int_0^\infty d\bC\,\bC\,J_2(q\bC)\,\chi(\bC),
\label{eq:FT}
\end{align}
with, for the Gaussian and for a sharp cutoff $\chi=\theta(r_0-|\bC|)$,
\begin{align}
T_{\rm gauss}(q)=\frac{2}{q^{2}}\Big(1-e^{-q^{2}r_0^{2}/4}\Big)
-\frac{r_0^{2}}{2}\,e^{-q^{2}r_0^{2}/4},
\qquad
T_{\theta}(q)=\frac{2}{q^{2}}\big(1-J_0(qr_0)\big)
-\frac{r_0}{q}J_1(qr_0).
\label{eq:T}
\end{align}
Both satisfy $T(q)\to2/q^{2}$ as $q\to\infty$ and $T(q)=\mathcal O(q^2)$
as $q\to0$.  With $q=k/\xi$,
\begin{align}
\boxed{\;
\sigma_{\rm CT}
=-2\pi P\,S_{ij}(\khat)
\left[\int_{0}^{1}d\xi\,\frac{1-\xi}{\xi}\,T\!\big(|k|/\xi\big)\right]
\Sigma_{ij}\;},
\qquad
\Sigma_{ij}=\int_{x,y,x'}
\frac{(x-y)_i\,(x'-y)_j}{(x-y)^{2}(x'-y)^{2}}
\big[U(y)-U(x')\big]^{ce'}\big[U(y)-U(x)\big]^{ce}\rho^{e'}\rho^{e}.
\label{eq:CTclosed}
\end{align}
Since $T(|k|/\xi)\to2\xi^{2}/k^{2}$ at $\xi\to0$, the $\xi$ integrand
behaves as $2\xi/k^{2}$ at the lower end: \emph{the counterterm carries
no $\ln V$} and the cutoff has already been removed in
\eqref{eq:CTclosed}.  In the distributional scheme $r_0\to\infty$ the
$\xi$ integral collapses to
$\int_0^1d\xi\,(1-\xi)/\xi\cdot2\xi^2/k^2=1/(3k^{2})$ and
\begin{align}
\sigma_{\rm CT}^{\,r_0\to\infty}
=-\frac{2\pi}{3k^{2}}\,P\,S_{ij}(\khat)\,\Sigma_{ij}.
\end{align}

\begin{remark}[Zero residue]
Had one defined the counterterm with the phase frozen to $1$ as well
(the textbook soft subtraction), it would vanish identically:
$\int d^2\bC\,\chi\,S_{ij}(\Chat)=0$ by property (ii) of
Sec.~\ref{sec:kernel}.  This is the subtraction-scheme restatement of
Sec.~\ref{sec:finite}: there is no rapidity divergence to subtract from
this term at $k\ne0$.  Eq.~\eqref{eq:CT} is a finite, $V$-independent
reorganization whose sole purpose is regularity of the integrand at
$y'=y$.
\end{remark}

\begin{remark}[Scheme dependence]
$r_0$ is a scheme parameter; its dependence cancels exactly between
$\sigma_{\rm reg}$ and $\sigma_{\rm CT}$.  A physical choice is the
typical dipole/target size.  The frozen factor $[U(y)-U(x')]$ in
\eqref{eq:CT} no longer vanishes at $y'=x'$, but the counterterm kernel
$(x'-y)_j/(x'-y)^2$ is regular there, so no new singularity is
introduced.  The $k\to0$ sensitivity of $\sigma$ is now explicit in the
$S_{ij}(\khat)/k^{2}$ of the counterterm, with a matching $\ln|k|$ in
$\sigma_{\rm reg}$.
\end{remark}

\section{Numerical verification}
\label{sec:numerics}

All claims were checked in a toy model that keeps every contested
structure and drops spectators: $y=0$ (so $\bC=y'$), fixed
$\bA=(1,0.2)$ and $x'=(-0.6,0.8)$ placed outside the integration domain
$|\bC|<\tfrac12$ (the $y'=x'$ point being separately protected in the
full expression), a smooth stand-in
$f(\bC)=e^{-(\bC-c_0)^2}$, $c_0=(0.13,0.21)$, for the Wilson-line and
source factors (nonzero and off-center at $\bC=0$ so that no accidental
symmetry hides anything), $k=(1,0)$, $\varepsilon=1/V$, and the exact
$F_\varepsilon$ of Appendix~\ref{app:F} (itself verified against direct
quadrature of the $\xi$ integral).  The scripts are in
\texttt{notes/checks/}.

\paragraph{(1) The slice diverges pointwise.}
$F_{1/V}(0)=\ln V-1+1/V$: $3.6152,\ 8.2104,\ 12.8155,\ 17.4207$ at
$V=10^{2},10^{4},10^{6},10^{8}$ --- steps of exactly $\ln100$.

\paragraph{(2) The integral converges (phase active, $k\ne0$).}
$I(V)=\int_{|\bC|<1/2}d^2\bC\,\mathcal K f\,F_{1/V}(k\cdot\bC)$:
\begin{center}
\begin{tabular}{lllll}
\toprule
$V$ & $10^{2}$ & $10^{4}$ & $10^{6}$ & $10^{8}$\\
\midrule
$I(V)$ & $-0.21637+0.01615i$ & $-0.21664+0.01616i$
       & $-0.21679+0.01616i$ & $-0.21679+0.01616i$\\
\bottomrule
\end{tabular}
\end{center}
While the integrand at $y'=y$ grew from $3.6$ to $17.4$, the integral
moved in the fourth decimal and froze.

\paragraph{(3) Remove the phase and the divergence is real.}
With $k=0$: $I=-0.0079,\ -0.0179,\ -0.0280,\ -0.0380$ at the same $V$
values --- linear in $\ln V$ ($-0.01006$ per factor $10^{2}$).  The
finiteness in (2) is the phase at work, not smallness of the kernel.

\paragraph{(4) Exact angular cancellation.}
$\oint d\theta\,\mathcal K(r\hat n)$ evaluates to $10^{-16}$ at every
radius tested ($r=0.2$ down to $2\times10^{-4}$), confirming
Eq.~\eqref{eq:loopzero} at all $r$, as Lemma~\ref{lem:loop} requires.

\paragraph{(5) The subtraction works.}
With $r_0=0.1$ and the frozen factor $f(0)$:
the Fourier identity \eqref{eq:FT}--\eqref{eq:T} was confirmed to eight
digits componentwise; the counterterm from the closed form
\eqref{eq:CTclosed}, $-0.0101750$, agrees with its direct
two-dimensional evaluation at $V=10^{8}$, $-0.0101754$, whose imaginary
part vanishes as predicted; the subtracted integral is flat,
$\sigma_{\rm reg}=-0.20663,\,-0.20649,\,-0.20662,\,-0.20662\;(+0.01616i)$
at $V=10^{2},10^{4},10^{6},10^{8}$; and
$\sigma_{\rm reg}+\sigma_{\rm CT}=-0.21679229+0.01616024i$ reproduces
the original $-0.21679230+0.01616024i$ to eight digits.  Finally, the
$y'$-density after the $\xi$ integration at $V=10^{12}$, along
$\bC=(\epsilon,0)$:
\begin{center}
\begin{tabular}{lll}
\toprule
$\epsilon$ & unsubtracted & subtracted\\
\midrule
$10^{-2}$ & $+2.100-1.042i$ & $+0.0283-0.0140i$\\
$10^{-3}$ & $+3.667-1.075i$ & $+0.00174-0.00051i$\\
$10^{-4}$ & $+5.248-1.079i$ & $+0.00020-0.00004i$\\
$10^{-5}$ & $+6.831-1.080i$ & $+0.000026-0.000004i$\\
\bottomrule
\end{tabular}
\end{center}
The unsubtracted density climbs like $\ln(1/\epsilon)$ toward the
$\ln V$ plateau; the subtracted one vanishes like
$\epsilon\ln(1/\epsilon)$.  The pole is gone pointwise.

\appendix

\section{Closed form of the longitudinal integral}
\label{app:F}

With $t=1/\xi$,
$F_\varepsilon(a)=\int_{1}^{1/\varepsilon}dt\,(t^{-1}-t^{-2})e^{-iat}$.
For $a=0$, $F_\varepsilon(0)=\ln(1/\varepsilon)-1+\varepsilon$.  For
$a\neq0$, writing $T=1/\varepsilon$,
\begin{align}
\int_1^T\frac{e^{-iat}}{t}\,dt
&=\mathrm{Ci}(|a|T)-\mathrm{Ci}(|a|)
-i\,\mathrm{sgn}(a)\big[\mathrm{Si}(|a|T)-\mathrm{Si}(|a|)\big]
\;\equiv\;g_1,\\
\int_1^T\frac{e^{-iat}}{t^{2}}\,dt
&=e^{-ia}-\frac{e^{-iaT}}{T}-i\,a\,g_1\;\equiv\;g_2
\qquad\text{(by parts)},\\
F_\varepsilon(a)&=g_1-g_2 .
\end{align}
The bounds \eqref{eq:Fbounds} follow from the asymptotics of
$\mathrm{Si}$, $\mathrm{Ci}$: the oscillation of $e^{-iat}$ truncates
the logarithmic growth of $g_1$ at $t\sim1/|a|$, so
$F_\varepsilon(a)=\ln\big(\min(1/|a|,\,1/\varepsilon)\big)+\mathcal O(1)$
uniformly, with $1/|a|$ decay for $|a|\gtrsim1$.

\end{document}
